\documentclass[10pt,a4paper]{article}
\usepackage{amsmath,amssymb,amsthm,graphicx}
\usepackage[margin=22mm]{geometry}
\usepackage{mathtools,graphicx,nicematrix,enumitem}
\usepackage[hidelinks]{hyperref}
\newtheorem{thm}{Theorem}[section]
\newtheorem{lem}[thm]{Lemma}
\theoremstyle{definition}
\newtheorem{dfn}[thm]{Definition}
\newcommand\myatop[2]{\genfrac{}{}{0pt}{}{#1}{#2}}
\newcommand{\myscale}{0.75}
\newcommand{\deff}[1]{\emph{#1}}
\DeclareMathOperator{\sgn}{sgn}
\DeclareMathOperator{\asym}{\mathbf{ASym}}
\DeclareMathOperator*{\pf}{Pf}
\newcommand{\boldk}{\boldsymbol{k}}
\title{Littlewood-type infinite-sum identity for modified Robbins polynomials (Theorem 1.1)}
\author{Ilse Fischer, Hans Höngesberg}
\date{doi:10.5802/alco.495; CC BY 4.0}
\begin{document}
\maketitle
\noindent\textit{Editorial note.} Mathematical spans are unchanged native published TeX. Counters restore original theorem/definition/lemma and global equation numbers. Full 43-entry bibliography recovered from the official article HTML, preserving numbers and original keys; math uses official data-tex. No external figure or current cited proof dependency in the selected chain.\par\medskip
\noindent\textit{Editorial scope.} One complete Theorem 1.1 with support; no mathematical repair, alternate-proof count or formal admission.\par\medskip
\noindent\textit{Allowed standard prerequisites (editorial).}\begin{itemize}
\item Pfaffian algebra. Source native 575–597, equation (6). Perfect-matching definition, skew-symmetric completion, simultaneous row/column operations and last-column expansion are retained.
\item Polynomial interpolation and divisibility; formal geometric series. Complete Lemma 4.1 proof native 616–723 and main proof 727–757. Standard techniques applied to the explicit current identities, which are proved in full; no external current core or computation is invoked.
\end{itemize}\medskip
\paragraph{Source and scope.} Ilse Fischer and Hans H\"ongesberg, \emph{A Littlewood-type identity for Robbins polynomials}, Algebraic Combinatorics 9(3) (2026), 665--700, \href{https://doi.org/10.5802/alco.495}{doi:10.5802/alco.495}. Copyright the authors, 2026; \href{https://creativecommons.org/licenses/by/4.0/}{CC BY 4.0}.
This extraction retains Theorem 1.1 as one counted result, Definition 2.1, the Pfaffian preliminaries, Lemma 4.1 with its full proof, and the complete first proof of Theorem 1.1. Original numbering and mathematical text are retained; the alternative proof sketch and subsequent consequences are not part of the counted package.
The author source is unchanged. Only a standard article wrapper and bibliography recovered from the same official article have been supplied. This is author-material preparation, not independent mathematical certification. The geometric-series interpretation is the one used in the retained author proof.

\setcounter{section}{1}
\setcounter{thm}{0}
\setcounter{equation}{1}

	\begin{thm}
		\label{thm:Robbins_Littlewood}
		Let $n$ be a positive integer. Then 
		\begin{multline}\label{eq:Robbins_Littlewood}
			\sum_{0 \le k_1 < \dots < k_n} R^{\ast}_{(k_{1},\ldots,k_{n})}(x_{1},\ldots,x_{n};1,1,w)\\
			=\prod_{i=1}^n \frac{1}{1-x_i}  \prod_{1 \le i < j \le n} \frac{x_i+x_j+w x_i x_j}{x_j-x_i}\\
			\times
			\pf_{\chi_{\mathrm{even}}(n) \le i < j \le n} \left(
			\begin{cases}
				1, & i=0,\\
				\frac{(x_j-x_i)(1+(1+w) x_i x_j)}{(x_i+x_j+w x_i x_j)(1-x_i x_j)},& i\ge 1,
			\end{cases} \right)
		\end{multline}
		where $\pf$ denotes the Pfaffian of an upper triangular array and $\chi_{\mathrm{even}}(n)$ equals $1$ if $n$ is even and $0$ otherwise.
	\end{thm}

\setcounter{section}{1}

	\section{Modified Robbins polynomials}
	\label{sec:robbins}
	
	To begin with, let us fix some notation. For any positive integer~$n$, we denote by $\mathfrak{S}_{n}$ the symmetric group of order~$n$. We define the \deff{antisymmetrizer} $\asym_{x_1,\ldots,x_n}$ to be 
	
	\begin{equation*}
		\asym_{x_1,\ldots,x_n} F(x_{1},\ldots,x_{n}) \coloneqq \sum_{\sigma \in \mathfrak{S}_{n}}\sgn(\sigma) F(x_{\sigma(1)},\ldots,x_{\sigma(n)}).
	\end{equation*}
	The antisymmetrizer is needed to define the modified Robbins polynomials, which are one of the central objects of this paper.
	
	\begin{dfn}
		Let $\boldk=(k_{1},\ldots,k_{n})$ be a sequence of integers. We define the \deff{modified Robbins polynomials} $R_{\boldk}^{\ast}(x_{1},\ldots,x_{n};u,v,w)$ to be
		\begin{equation*}
			\frac{\asym_{x_{1},\ldots,x_{n}}\left[\prod\limits_{1\le i<j\le n}(ux_{i}x_{j}+v+wx_{i})\prod\limits_{i=1}^{n}x_{i}^{k_{i}}\right]}{\prod\limits_{1\le i<j\le n}(x_{j}-x_{i})}.
		\end{equation*}
	\end{dfn}

\setcounter{section}{3}
\section{Proof of Theorem~\ref{thm:Robbins_Littlewood}}
\label{sec:mainproof}
\setcounter{equation}{5}

	The Littlewood-type identity for modified Robbins polynomials  involves the notion of Pfaffians which we recall in the following before proving Theorem~\ref{thm:Robbins_Littlewood}:
	
	For every triangular array $A=(a_{i,j})_{1 \le i < j \le 2n}$, the Pfaffian $\pf(A)$ is defined as
	\begin{equation*}
		\sum_{\{(i_1,j_1),\dots,(i_n,j_n)\}} \sgn(i_1 j_1 \dots i_n j_n) a_{i_1,j_1} \cdots a_{i_n,j_n},
	\end{equation*}
	where we sum over all perfect matchings $\{(i_1,j_1),\dots,(i_n,j_n)\}$ of $\{1,\dots,2n\}$ such that $i_1 < \dots < i_n$ and $i_k < j_k$ for all $1 \le k \le n$.
	
	Every triangular array $A=(a_{i,j})_{1 \le i < j \le 2n}$ is the upper triangular part of a uniquely determined $2n \times 2n$ skew-symmetric matrix $M$, which we call the \emph{skew-symmetric completion} of $A$. A matrix $M=(m_{i,j})_{1 \le i , j \le 2n}$ is skew-symmetric if $m_{i,j} = - m_{j,i}$ for all $1 \le i,j \le 2n$; in particular, all diagonal entries are $0$. It is well known that $\pf(A)^2=\det(M)$. On that account, we will repeatedly change perspectives when dealing with Pfaffians as we will consider either the triangular array or the skew-symmetric matrix that underlies the Pfaffian.
	
	We recall the main methods for manipulating Pfaffians used in the present paper:
	\begin{itemize}
		\item Regarding elementary row or column operations on the underlying  matrix, we perform corresponding row and column operations at the same time. Thus, the resulting matrix remains skew-symmetric. In particular, simultaneously swapping corresponding rows and columns multiplies the Pfaffian by $-1$.
		\item Expanding the Pfaffian along the last column yields
		\begin{equation*}
			\pf_{1 \le i < j \le 2n} (a_{i,j}) = \sum_{k=1}^{2n-1} (-1)^{k+1} a_{k,2n} \pf_{\myatop{1 \le i < j \le 2n}{i,j \neq k}} (a_{i,j}).
		\end{equation*}
		\item For all $c_1,\dots,c_{2n}$, it holds that
		\begin{equation}\label{eq:linearity_pfaffian}
			\pf_{1 \le i < j \le 2n} (c_i c_j a_{i,j}) = \prod_{i=1}^{2n} c_i \pf_{1 \le i < j \le 2n} (a_{i,j}).
		\end{equation}
		Thus, for a given~$i$, setting $c_i=c$ and all other $c_1,\dots,c_{i-1},c_{i+1},\dots,c_{2n}$ equal to $1$ corresponds to multiplying the $i$th row and the $i$th column of the underlying matrix by $c$.
	\end{itemize}
	
	In preparation for the proof of Theorem~\ref{thm:Robbins_Littlewood}, we show the following lemma which will be used for certain row and column operations:

	\begin{lem}\label{lem:MatrixOperationEven}
		The following identities hold:
		\begin{enumerate}[label=(\roman*)]
			\item\label{eq:MatrixOperationEven2} For each positive integer~$n$, it holds
			\begin{equation*}
				\sum_{j=1}^{n} \left(\frac{1+x_j}{x_j} \prod_{\myatop{1 \le p \le n}{p \neq j}} \frac{1-x_j x_p}{x_j -x_p} \right) + (-1)^n \prod_{k=1}^{n} \frac{1}{x_k} = 1.
			\end{equation*}
			\item\label{eq:MatrixOperationEven1} For each positive integer~$n$ and $i$ satisfying $1 \le i \le n$, it holds
			\begin{multline*}
				(-1)^n \sum_{j=1}^{n} \left( \frac{(x_j-x_i)(1+(1+w) x_i x_j)}{(x_i+x_j+w x_i x_j)(1-x_i x_j)} \cdot \frac{1+x_j}{x_j} \prod_{\myatop{1 \le p \le n}{p \neq j}} \frac{1-x_j x_p}{x_j -x_p} \right)\\
				+ \frac{1-x_i}{x_i} \prod\limits_{\myatop{1 \le p \le n}{p \neq i}} \frac{1+ w x_i + x_i x_p}{x_i + x_p + w x_i x_p} = (-1)^n+\prod_{k=1}^{n} \frac{1}{x_k}.
			\end{multline*}
		\end{enumerate}
	\end{lem}
	
	\begin{proof}
		Let us first prove~\ref{eq:MatrixOperationEven2}. After some straightforward manipulation, the expression in~\ref{eq:MatrixOperationEven2} becomes the following polynomial identity:
		\begin{multline}\label{eq:MatrixOperationEven2_1}
			\sum_{j=1}^{n} \left((-1)^j (1+x_j) \prod_{\myatop{1 \le p \le n}{p \neq j}} x_p(1-x_j x_p) \prod_{\myatop{1 \le p < q\le n}{p,q \neq j}} (x_q - x_p) \right)\\
			=  \left( -1 + (-1)^n \prod_{p=1}^{n} x_p \right)  \prod_{1 \le p < q \le n} (x_q - x_p).
		\end{multline}
		Both sides are multivariate polynomials in $x_1,\dots,x_n$ with highest degree at most $n$ in each variable. It is not difficult to see that both sides vanish if we set $x_s=x_t$ for any $s \neq t$. Furthermore, if $x_s=0$, both sides evaluate to 
		\begin{equation*}
			(-1)^s  \prod_{\myatop{1 \le p \le n}{p \neq s}} x_p \prod_{\myatop{1 \le p < q \le n}{p,q \neq s}} (x_q - x_p),
		\end{equation*}
		resulting in $n$ coinciding evaluations for each variable. Thus, both sides are equal up to a multiplicative factor independent of $x_1,\dots,x_n$. Since the coefficient of $x_n^n x_{n-1}^{n-1} \cdots x_1$, that is, the leading coefficient in the lexicographic order for $x_n > \dots > x_1$, equals $(-1)^n$ on both sides, we have shown~\eqref{eq:MatrixOperationEven2_1}.
		
		Regarding~\ref{eq:MatrixOperationEven1}, we transform the identity into
		\begin{multline}\label{eq:MatrixOperationEven1_1}
			\sum_{j=1}^{n} \left( \vphantom{ \prod_{\myatop{1 \le p < q \le n}{p,q \neq j}}}  (-1)^{j} (1+(1+w) x_i x_j) (1+x_j) (x_j-x_i) \right.\\
			\left. \times \prod_{\myatop{1 \le p \le n}{p \neq j}} x_p \prod_{\myatop{1 \le p \le n}{p \neq i,j}} (1-x_j x_p)  (x_i+x_p+w x_i x_p) \prod_{\myatop{1 \le p < q \le n}{p,q \neq j}} (x_q - x_p) \right) \\
			+ (1-x_i)  \prod\limits_{\myatop{1 \le p \le n}{p \neq i}} x_p (1+ w x_i + x_i x_p) \prod_{1 \le p <q \le n} (x_q -x_p)\\
			= \left(1+ (-1)^n \prod_{k=1}^{n} x_k \right) \prod_{\myatop{1 \le p \le n}{p \neq i}} (x_i+x_p+w x_i x_p) \prod_{1 \le p < q \le n } (x_q -x_p),
		\end{multline}
		both sides of which are multivariate polynomials in $x_1,\dots,x_n$; the highest degree of $x_s$ is at most $n+1$ if $s \neq i$ and $2n-1$ otherwise.
		
		A straightforward computation shows that both sides vanish if we set $x_s=x_t$ for any $s \neq t$. Setting $x_s=0$ yields
		\begin{equation*}
			(-1)^{s-1} x_i^2 \prod_{\myatop{1 \le p \le n}{p \neq i,s}} x_p (x_i+x_p+w x_i x_p) \prod_{\myatop{1 \le p < q \le n}{p,q \neq s}} (x_q - x_p)
		\end{equation*}
		on both sides if $s \neq i$; otherwise, we obtain
		\begin{equation*}
			(-1)^{i-1}  \prod_{\myatop{1 \le p \le n}{p \neq i}} x_p^2 \prod_{\myatop{1 \le p < q \le n}{p,q \neq i}} (x_q - x_p).
		\end{equation*}
		
		Next, we consider the evaluation for $x_s=-\frac{x_i}{1+w x_i}$ if $s \neq i$. The right-hand side of~\eqref{eq:MatrixOperationEven1_1} clearly vanishes. The left-hand side, however, becomes
		\begin{multline*}
			\frac{(-1)^{s}}{(1+w x_i)^{n+1}} x_i (x_i-1)(2+ w x_i)(1-x_i+w x_i)(1+x_i+w x_i)\\
			\times \prod_{\myatop{1 \le p \le n}{p \neq s}} x_p \prod_{\myatop{1 \le p \le n}{p \neq i,s}} (1+w x_i+ x_i x_p)  (x_i+x_p+w x_i x_p) \prod_{\myatop{1 \le p < q \le n}{p,q \neq s}} (x_q - x_p)\\
			- \frac{(-1)^s}{(1+w x_i)^{n+1}} x_i (x_i-1) (1-x_i+w x_i)(1+x_i + w x_i)\\
			\times\prod_{\myatop{1 \le p \le n}{p \neq s}}  (x_i+x_p + w x_i x_p)  \prod_{\myatop{1 \le p \le n}{p \neq i,s}} x_p (1+ w x_i + x_i x_p) \prod_{\myatop{1 \le p <q \le n}{p,q \neq s}} (x_q -x_p).
		\end{multline*}
		After omitting common factors, we are left with $(2+ w x_i) x_i -  (x_i+x_i + w x_i^2)$, which equals zero. 
		
		Next, let $x_i=-\frac{x_s}{1+w x_s}$ for any $s \neq i$. It is again obvious that the right-hand side of~\eqref{eq:MatrixOperationEven1_1} vanishes. We obtain for the left-hand side
		\begin{multline*}
			\frac{(-1)^{s+i-\left[s<i\right]}}{(1+w x_s)^{2n-1}} x_s^2 (1-x_s) (1+x_s) (2+ w x_s)(1+x_s+w x_s)\\
			\times  \prod_{\myatop{1 \le p \le n}{p \neq i,s}}x_p (1-x_s x_p)  (x_p - x_s)  (x_s+x_p+w x_s x_p) \prod_{\myatop{1 \le p < q \le n}{p,q \neq i,s}} (x_q - x_p)\\
			+ \frac{(-1)^{i-1}}{(1+w x_s)^{2n-1}}  (1+x_s + w x_s)\prod_{\myatop{1 \le p \le n}{p \neq i}} x_p (1-x_s x_p) (x_s+x_p+w x_s x_p) \\
			\times \prod_{\myatop{1 \le p <q \le n}{p,q \neq i}} (x_q -x_p),
		\end{multline*}
		where $\left[\boldsymbol{\cdot}\right]$ denotes the \emph{Iverson bracket}: For any proposition~$P$, $\left[P\right]$ evaluates to $1$ if $P$ is true and to $0$ if $P$ is false. It is again not difficult to see that the previous expression is zero.
		
		In total, we have now $n+1$ matching evaluations for $x_s$ for any $s\neq i$ and $2n-1$ evaluations for $x_i$. As a consequence, we know that the difference between the left-hand side and the right-hand side of~\eqref{eq:MatrixOperationEven1_1} is the product of
		\begin{equation}\label{eq:differencepoly}
			\prod_{p=1}^{n} x_p \prod_{1 \le p < q \le n} (x_q-x_p) \prod_{\myatop{1 \le p \le n}{p \neq i}} (x_i + x_p + w x_i x_p)^2
		\end{equation}
		and a rational expression in $w$. In fact, we want to show that this factor equals zero. To this end, we consider the specializations $(x_1,\dots,x_n)=(x_1,\frac{1}{x_1},x_3,\frac{1}{x_3},\dots,x_{n-1},\frac{1}{x_{n-1}})$ if $n$ is even and $(x_1,\dots,x_n)=(x_1,\frac{1}{x_1},x_3,\frac{1}{x_3},\dots,x_{n-2},\frac{1}{x_{n-2}},-1)$ if $n$ is odd for which~\eqref{eq:differencepoly} does not vanish. Given these specializations, we will show that the left-hand side of~\eqref{eq:MatrixOperationEven1_1} coincides with the right-hand side, which completes the proof of~\eqref{eq:MatrixOperationEven1_1}.
		
		First, we consider the left-hand side of~\eqref{eq:MatrixOperationEven1_1}. Due to the factor $\prod_{\myatop{1 \le p \le n}{p \neq i,j}} (1-x_j x_p)$, we see that the first sum vanishes unless $j=i+1$ if $i$ is odd and $j=i-1$ if $i$ is even. Let us assume the latter. Thus, it suffices to only keep the $(i-1)$th summand in the first sum. After some manipulation, the left-hand side of~\eqref{eq:MatrixOperationEven1_1} reduces to
		\begin{multline*}
			x_i (1+(1+w) x_{i-1} x_i) (1+x_{i-1}) (x_i-x_{i-1})\\
			\times \prod_{\myatop{1 \le p \le n}{p \neq i-1,i}} x_p (1-x_{i-1} x_p)  (x_i+x_p+w x_i x_p) (x_p - x_i) \prod_{\myatop{1 \le p < q \le n}{p,q \neq i-1,i}} (x_q - x_p) \\
			+ x_{i-1} (1-x_i) (1+w x_i+x_{i-1} x_i) (x_i-x_{i-1})\\
			\times\prod\limits_{\myatop{1 \le p \le n}{p \neq i-1,i}} x_p (1+ w x_i + x_i x_p) (x_p-x_i) (x_p-x_{i-1}) \prod_{\myatop{1 \le p <q \le n}{p,q \neq i-1,i}} (x_q -x_p).
		\end{multline*}
		
		The right-hand side of~\eqref{eq:MatrixOperationEven1_1}, however, can be recasted as
		\begin{multline*}
			(x_{i-1}+x_i+w x_{i-1} x_i) (x_i-x_{i-1})\\
			\times \left(1+ (-1)^n  \prod_{k=1}^{n} x_k \right)\prod\limits_{\myatop{1 \le p \le n}{p \neq i-1,i}} (x_i+ x_p +w  x_i x_p) (x_p-x_i) (x_p-x_{i-1}) \prod_{\myatop{1 \le p <q \le n}{p,q \neq i-1,i}} (x_q -x_p).
		\end{multline*}
		After cancelling out common factors in the two previous expressions, it remains to show that
		\begin{multline*}
			x_i (1+(1+w) x_{i-1} x_i) (1+x_{i-1}) \prod_{\myatop{1 \le p \le n}{p \neq i-1,i}} x_p (1-x_{i-1} x_p)  (x_i+x_p+w x_i x_p) \\
			+ x_{i-1} (1-x_i) (1+w x_i+x_{i-1} x_i) \prod\limits_{\myatop{1 \le p \le n}{p \neq i-1,i}} x_p (1+ w x_i + x_i x_p) (x_p-x_{i-1}) 
		\end{multline*}
		equals
		\begin{equation*}
			(x_{i-1}+x_i+w x_{i-1} x_i) \left(1+ (-1)^n \prod_{k=1}^{n} x_k \right)\prod\limits_{\myatop{1 \le p \le n}{p \neq i-1,i}} (x_i+ x_p +w  x_i x_p) (x_p-x_{i-1}) 
		\end{equation*}
		for the specializations under consideration. For this purpose, let $k \neq i$ be even and consider the products
		\begin{itemize}
			\item $x_{k-1} (1-x_{i-1} x_{k-1})  (x_i+x_{k-1}+w x_i x_{k-1}) \cdot x_{k} (1-x_{i-1} x_{k})  (x_i+x_{k}+w x_i x_{k})$,
			\item $x_{k-1} (1+ w x_i + x_i x_{k-1}) (x_{k-1}-x_{i-1}) \cdot x_{k} (1+ w x_i + x_i x_{k}) (x_{k}-x_{i-1})$,
			\item $(x_i+ x_{k-1} +w  x_i x_{k-1}) (x_{k-1}-x_{i-1}) \cdot (x_i+ x_{k} +w  x_i x_{k}) (x_{k}-x_{i-1})$,
		\end{itemize}
		all of which evaluate to
		\begin{equation*}
			\frac{1}{x_{k-1}^2 x_i^2} (x_{k-1}-x_{i-1})(1-x_{i-1} x_{k-1})(w+x_{i-1}+x_{k-1})(1+w x_{k-1}+x_{i-1} x_{k-1})
		\end{equation*}
		when setting $x_i \mapsto \frac{1}{x_{i-1}}$ and $x_k \mapsto \frac{1}{x_{k-1}}$. In addition, in the case of $n$ being odd, one can easily check that all of the following expressions
		\begin{itemize}
			\item $x_{n} (1-x_{i-1} x_{n})  (x_i+x_{n}+w x_i x_{n})$,
			\item $x_{n} (1+ w x_i + x_i x_{n}) (x_{n}-x_{i-1})$,
			\item $(x_i+ x_{n} +w  x_i x_{n}) (x_{n}-x_{i-1})$
		\end{itemize}
		are equal to $(1+x_{i-1})(1 - x_i + w x_i)$ if setting $x_n=-1$.
		
		Moreover, we see that
		\begin{equation*}
			x_i (1+(1+w) x_{i-1} x_i) (1+x_{i-1}) + x_{i-1} (1-x_i) (1+w x_i+x_{i-1} x_i)  
		\end{equation*}
		and
		\begin{equation*}
			(x_{i-1}+x_i+w x_{i-1} x_i) \left(1+ (-1)^n \prod_{k=1}^{n} x_k \right)
		\end{equation*}
		both specialize to
		\begin{equation*}
			\frac{2}{x_{i-1}} (1+ w x_{i-1}+ x_{i-1}^2).
		\end{equation*}
		This finally concludes the proof of~\eqref{eq:MatrixOperationEven1_1} for even $i$. An analogous argument holds in the case of $i$ being odd.
	\end{proof}
	
	We are now ready to proceed with the main proof of Theorem~\ref{thm:Robbins_Littlewood}.
	
	\begin{proof}[First proof of Theorem~\ref{thm:Robbins_Littlewood}]
		We prove the statement by induction on $n$, showing that both sides of~\eqref{eq:Robbins_Littlewood} satisfy the same recurrence relation.
		
		To start with, recall that $R_{\boldk}^{\ast}(x_{1},\ldots,x_{n};1,1,w)$ is given by
		\begin{equation*}
			\frac{\asym_{x_{1},\ldots,x_{n}}\left[\prod_{1\le i<j\le n}(x_{i}x_{j}+1+wx_{i})\prod_{i=1}^{n}x_{i}^{k_{i}}\right]}{\prod_{1\le i<j\le n}(x_{j}-x_{i})}.
		\end{equation*}
		Through geometric series evaluation, it follows that
		\begin{equation*}
			F(x_1,\ldots,x_n) \coloneqq \sum_{0 \le k_1 < \dots < k_n} R^{\ast}_{(k_{1},\ldots,k_{n})}(x_{1},\ldots,x_{n};1,1,w)
		\end{equation*} equals
		\begin{equation}\label{eq:asym_expression} 
			\frac{\asym_{x_{1},\ldots,x_{n}}\left[\prod_{1\le i<j\le n}(x_{i}x_{j}+1+wx_{i})\prod_{i=1}^{n}x_{i}^{i-1}(1-\prod_{j=i}^n x_j)^{-1}\right]}{\prod_{1\le i<j\le n}(x_{j}-x_{i})} 
		\end{equation}
		since 	
		\begin{multline*}
			\sum_{0 \le k_1 < \dots < k_n} \prod_{i=1}^{n} x_i^{k_i} =\prod_{i=1}^{n} x_i^{i-1} \sum_{0 \le k_1 \le \dots \le k_n} \prod_{i=1}^{n} x_i^{k_i} \\
			= \prod_{i=1}^{n} x_i^{i-1} \sum_{0 \le k_1 \le \dots \le k_n} (x_1 \dots x_n)^{k_1} (x_2 \dots x_n)^{k_2-k_1} \cdots x_n^{k_n-k_{n-1}}\\
			= \prod_{i=1}^{n} x_i^{i-1} \sum_{k_1=0}^{\infty} (x_1 \dots x_n)^{k_1} \sum_{k_2=0}^{\infty} (x_2 \dots x_n)^{k_2} \cdots  \sum_{k_n=0}^{\infty} x_n^{k_n} = \prod_{i=1}^{n} \frac{x_i^{i-1}}{1-\prod_{j=i}^{n} x_j}.
		\end{multline*}
		Thus, we obtain the recurrence relation
		\begin{equation}\label{eq:rec_Littlewood}
			F(x_1,\ldots,x_n) = 
			\left(1-\prod_{i=1}^n x_i \right)^{-1}  \sum\limits_{k=1}^n \prod\limits_{\myatop{1 \le i \le n}{i \neq k}} \frac{(x_i x_k+1+w x_k) x_i}{x_i-x_k}
			F(x_1,\ldots,\widehat{x_k},\ldots,x_n)
		\end{equation}
		with initial condition
		\begin{equation*}
			F(x_1)=\frac{1}{1-x_1},
		\end{equation*}
		where $x_1,\ldots,\widehat{x_k},\ldots,x_n$ denotes the omission of $x_k$ in $x_1,\ldots,x_n$. 
		
		\NiceMatrixOptions{cell-space-limits=5pt}
		
		If $n=1$, then both sides of~\eqref{eq:Robbins_Littlewood} clearly evaluate to $1/(1-x_1)$. Now consider the induction step from $n-1$ to $n$. We make a case distinction regarding the parity of $n$.
		
		Let us first assume that $n$ is even. Using the induction hypothesis and the recursion~\eqref{eq:rec_Littlewood}, it follows that the right-hand side of~\eqref{eq:Robbins_Littlewood} equals
		\begin{multline*}
			\left(1-\prod_{i=1}^n x_i\right)^{-1} \sum\limits_{k=1}^n \prod\limits_{\myatop{1 \le i \le n}{i \neq k}} \frac{(x_i x_k+w x_k + 1) x_i}{(x_i-x_k)(1-x_i)}\\
			\times \prod_{\myatop{1 \le i < j \le n}{i,j \neq k}} \frac{x_i+x_j+w x_i x_j}{x_j-x_i} 
			\left. \pf_{\myatop{1 \le i < j \le n+1}{i,j \neq k}} \left( \frac{(x_j-x_i)(1+(1+w) x_i x_j)}{(x_i+x_j+w x_i x_j)(1-x_i x_j)} \right)\right|_{x_{n+1}=1},
		\end{multline*}
		which we rewrite as
		\begin{multline*}
			\left(1-\prod_{i=1}^n x_i\right)^{-1} \prod\limits_{1 \le i < j \le n} \frac{x_i + x_j + w x_i x_j}{x_j-x_i} \sum\limits_{k=1}^n (-1)^{k+1} \left( \prod\limits_{\myatop{1 \le i \le n}{i \neq k}} \frac{1+ w x_k + x_k x_i}{x_k + x_i + w x_k x_i} \right)\\
			\times \prod\limits_{\myatop{1 \le i \le n}{i \neq k}} \frac{x_i}{1-x_i} 
			\left. \pf_{\myatop{1 \le i < j \le n+1}{i,j \neq k}} \left( \frac{(x_j-x_i)(1+(1+w) x_i x_j)}{(x_i+x_j+w x_i x_j)(1-x_i x_j)} \right)\right|_{x_{n+1}=1}.
		\end{multline*}
		Interpreting the sum as an expansion along the last column yields	
		\begin{multline*}
			\left(1-\prod_{i=1}^n x_i\right)^{-1} \prod\limits_{1 \le i < j \le n} \frac{x_i + x_j + w x_i x_j}{x_j-x_i}\\
			\times \pf_{1 \le i < j \le n+2}
			{\scalebox{\myscale}{$\begin{pmatrix}
						\begin{NiceArray}{@{}c|c|c@{}}[corners]
							\left( \frac{x_i}{1-x_i}  \frac{x_j}{1-x_j} \frac{(x_j-x_i)(1+(1+w) x_i x_j)}{(x_i+x_j+w x_i x_j)(1-x_i x_j)} \right)_{1 \le i < j \le n} & \Block{2-1}{\left( \frac{x_i}{1-x_i} \right)_{1 \le i \le n}} & \Block{2-1}{\left( \prod\limits_{\myatop{1 \le p \le n}{p \neq i}} \frac{1+ w x_i + x_i x_p}{x_i + x_p + w x_i x_p} \right)_{1\le i \le n}} \\ \cline{1}
							& & \\ \cline{2-3}
							& & 0
						\end{NiceArray}
					\end{pmatrix}$}}.	
		\end{multline*}
		We extract $\prod_{i=1}^{n} \frac{x_i}{1-x_i}$ from the Pfaffian using~\eqref{eq:linearity_pfaffian} and split it up into two summands as follows:	
		\begin{multline*}
			\pf_{1 \le i < j \le n+2} {\scalebox{\myscale}{$\begin{pmatrix}
						\begin{NiceArray}{@{}c|c|c@{}}[corners]
							\left( \frac{(x_j-x_i)(1+(1+w) x_i x_j)}{(x_i+x_j+w x_i x_j)(1-x_i x_j)} \right)_{1 \le i < j \le n} & \Block{2-1}{\left( 1 \right)_{1 \le i \le n}} & \Block{2-1}{\left( \frac{1-x_i}{x_i}  \left(\prod\limits_{\myatop{1 \le p \le n}{p \neq i}} \frac{1+ w x_i + x_i x_p}{x_i + x_p + w x_i x_p} \right) \right)_{1 \le i \le n}} \\ \cline{1}
							& & \\ \cline{2-3}
							& & -\prod_{k=1}^{n} \frac{1}{x_k}
						\end{NiceArray}
					\end{pmatrix}$}}\\
			+\prod\limits_{i=1}^{n} \frac{1}{x_i} \pf_{1 \le i < j \le n} \left(  \frac{(x_j-x_i)(1+(1+w) x_i x_j)}{(x_i+x_j+w x_i x_j)(1-x_i x_j)} \right).
		\end{multline*}
		This implies that the proof of the even case is complete if we show that	
		\begin{multline}\label{eq:2Pfaffians}
			\pf_{1 \le i < j \le n+2}  {\scalebox{\myscale}{$\begin{pmatrix}
						\begin{NiceArray}{@{}c|c|c@{}}[corners]
							\left( \frac{(x_j-x_i)(1+(1+w) x_i x_j)}{(x_i+x_j+w x_i x_j)(1-x_i x_j)} \right)_{1 \le i < j \le n} & \Block{2-1}{\left( -1 \right)_{1 \le i \le n}} & \Block{2-1}{\left( \frac{1-x_i}{x_i} \left( \prod\limits_{\myatop{1 \le p \le n}{p \neq i}} \frac{1+ w x_i + x_i x_p}{x_i + x_p + w x_i x_p} \right) \right)_{1 \le i \le n}} \\ \cline{1}
							& & \\ \cline{2-3}
							& & \prod_{k=1}^{n} \frac{1}{x_k}
						\end{NiceArray}
					\end{pmatrix}$}}\\
			= \pf_{1 \le i < j \le n} \left(  \frac{(x_j-x_i)(1+(1+w) x_i x_j)}{(x_i+x_j+w x_i x_j)(1-x_i x_j)} \right),
		\end{multline}
		where we incorporated a negative sign into the Pfaffian on the left-hand side by multiplying the $(n+1)$th row and $(n+1)$th column by $-1$ using~\eqref{eq:linearity_pfaffian}.
		
		We perform the following elementary row and column operations on the matrix underlying the Pfaffian in the left-hand side of~\eqref{eq:2Pfaffians}: For every $j\in\{1,\dots,n\}$, we add the $j$th column times
		\begin{equation*}
			\frac{1+x_j}{x_j} \prod_{\myatop{1 \le p \le n}{p \neq j}} \frac{1-x_j x_p}{x_j -x_p}
		\end{equation*}
		to the last column; we then proceed analogously for the rows. Due to Lemma~\ref{lem:MatrixOperationEven}, this results in the left-hand side of~\eqref{eq:2Pfaffians} being equal to
		
		\begin{equation*}\label{eq:2Pfaffians_1}
			\pf_{1 \le i < j \le n+2}   {\scalebox{\myscale}{$\begin{pmatrix}
						\begin{NiceArray}{@{}c|c|c@{}}[corners]
							\left( \frac{(x_j-x_i)(1+(1+w) x_i x_j)}{(x_i+x_j+w x_i x_j)(1-x_i x_j)} \right)_{1 \le i < j \le n} & \Block{2-1}{\left( -1 \right)_{1 \le i \le n}} & \Block{2-1}{\left(1+\prod_{k=1}^{n} \frac{1}{x_k} \right)_{1 \le i \le n}} \\ \cline{1}
							& & \\ \cline{2-3}
							& & 1
						\end{NiceArray}
					\end{pmatrix}$}}.
		\end{equation*}
		
		We further manipulate the underlying matrix as follows: We multiply the $(n+1)$th column by $1+\prod_{i=1}^{n} \frac{1}{x_i}$ and add it to the $(n+2)$th column; then we proceed similarly with the last two rows. Thus, we obtain
		\begin{equation*}
			\pf_{1 \le i < j \le n+2}  {\scalebox{\myscale}{$\begin{pmatrix}
						\begin{NiceArray}{@{}c|c|c@{}}[corners]
							\left( \frac{(x_j-x_i)(1+(1+w) x_i x_j)}{(x_i+x_j+w x_i x_j)(1-x_i x_j)} \right)_{1 \le i < j \le n} & \Block{2-1}{\left( -1 \right)_{1 \le i \le n}} & \Block{2-1}{\left(0 \right)_{1 \le i \le n}} \\ \cline{1}
							& & \\ \cline{2-3}
							& & 1
						\end{NiceArray}
					\end{pmatrix}$}}.
		\end{equation*}
		Evaluating this Pfaffian along the last column, finally yields the right-hand side of~\eqref{eq:2Pfaffians}.
		
		On the other hand, let us assume that $n$ is odd. Similar as before, applying the recursion~\eqref{eq:rec_Littlewood} to the induction hypothesis and some manipulations imply the right-hand side of~\eqref{eq:Robbins_Littlewood} being equal to
		\begin{multline*}
			\left(1-\prod_{i=1}^n x_i \right)^{-1} \prod\limits_{1 \le i < j \le n} \frac{x_i + x_j + w x_i x_j}{x_j-x_i} \sum\limits_{k=1}^n (-1)^{k+1} \left( \prod\limits_{\myatop{1 \le i \le n}{i \neq k}} \frac{1+ w x_k + x_k x_i}{x_k + x_i + w x_k x_i} \right)\\
			\times \prod\limits_{\myatop{1 \le i \le n}{i \neq k}} \frac{x_i}{1-x_i} 
			\pf_{\myatop{1 \le i < j \le n}{i,j \neq k}} \left( \frac{(x_j-x_i)(1+(1+w) x_i x_j)}{(x_i+x_j+w x_i x_j)(1-x_i x_j)} \right).
		\end{multline*}
		We multiply the factor $\prod_{\myatop{1 \le i \le n}{i \neq k}} \frac{x_i}{1-x_i} $ into the Pfaffian using~\eqref{eq:linearity_pfaffian} and then interpret the sum as an expansion along the last column, which yields
		\begin{multline*}
			\left(1-\prod_{i=1}^n x_i \right)^{-1} \prod\limits_{1 \le i < j \le n} \frac{x_i + x_j + w x_i x_j}{x_j-x_i}\\
			\times \pf_{1 \le i < j \le n+1}  {\scalebox{\myscale}{$\begin{pmatrix}
						\begin{NiceArray}{@{}c|c@{}}[corners]
							\left( \frac{x_i}{1-x_i}  \frac{x_j}{1-x_j} \frac{(x_j-x_i)(1+(1+w) x_i x_j)}{(x_i+x_j+w x_i x_j)(1-x_i x_j)} \right)_{1 \le i < j \le n} & \Block{2-1}{\left( \prod\limits_{\myatop{1 \le p \le n}{p \neq i}} \frac{1+ w x_i + x_i x_p}{x_i + x_p + w x_i x_p} \right)_{1 \le i \le n}} \\ \cline{1}
							&
						\end{NiceArray}
					\end{pmatrix}$}}.
		\end{multline*}
		Next, we extract again  the product $\prod_{i=1}^{n} \frac{x_i}{1-x_i}$ from the Pfaffian using~\eqref{eq:linearity_pfaffian}
		and recast it by enlarging the dimension of the underlying array as	
		\begin{equation*}
			\pf\limits_{1 \le i < j \le n+3}{\scalebox{\myscale}{$\begin{pmatrix}
						\begin{NiceArray}{@{}c|c|c|c@{}}[corners]
							\left( \frac{(x_j-x_i)(1+(1+w) x_i x_j)}{(x_i+x_j+w x_i x_j)(1-x_i x_j)} \right)_{1 \le i < j \le n} & \Block{2-1}{\left( 1 \right)_{1 \le i \le n}} & \Block{2-1}{\left( \frac{1-x_i}{x_i} \left( \prod\limits_{\myatop{1 \le p \le n}{p \neq i}} \frac{1+ w x_i + x_i x_p}{x_i + x_p + w x_i x_p} \right) \right)_{1 \le i \le n}}  & \Block{2-1}{(0)_{1 \le i \le n}} \\ \cline{1}
							& & & \\ \cline{2-4}
							& & 0 & -1 \\ \cline{3-4}
							& & & 0
						\end{NiceArray}
					\end{pmatrix}$}};
		\end{equation*}
		we then split up the Pfaffian into the following sum:	
		\begin{multline*}
			\pf_{1 \le i < j \le n+3}{\scalebox{\myscale}{$\begin{pmatrix}
						\begin{NiceArray}{@{}c|c|c|c@{}}[corners]
							\left( \frac{(x_j-x_i)(1+(1+w) x_i x_j)}{(x_i+x_j+w x_i x_j)(1-x_i x_j)} \right)_{1 \le i < j \le n} & \Block{2-1}{\left( 1 \right)_{1 \le i \le n}} & \Block{2-1}{ \left( \frac{1-x_i}{x_i} \left( \prod\limits_{\myatop{1 \le p \le n}{p \neq i}} \frac{1+ w x_i + x_i x_p}{x_i + x_p + w x_i x_p} \right) \right)_{1 \le i \le n}}  & \Block{2-1}{(0)_{1 \le i \le n}} \\ \cline{1}
							& & & \\ \cline{2-4}
							& & 0 & -1 \\ \cline{3-4}
							& & & -\prod\limits_{k=1}^{n} \frac{1}{x_k}
						\end{NiceArray}
					\end{pmatrix}$}}\\
			+\prod\limits_{i=1}^{n} \frac{1}{x_i} \pf_{1 \le i < j \le n+1}  {\scalebox{\myscale}{$\begin{pmatrix}
						\begin{NiceArray}{@{}c|c@{}}[corners]
							\left( \frac{(x_j-x_i)(1+(1+w) x_i x_j)}{(x_i+x_j+w x_i x_j)(1-x_i x_j)} \right)_{1 \le i < j \le n} & \Block{2-1}{\left(1 \right)_{1 \le i \le n}} \\ \cline{1}
							&
						\end{NiceArray}
					\end{pmatrix}$}}.
		\end{multline*}
		After putting all together and some minor manipulations, it remains to show that	
		\begin{multline}\label{eq:2Pfaffians_Odd}\\
			\pf_{1 \le i < j \le n+3}{\scalebox{\myscale}{$\begin{pmatrix}
						\begin{NiceArray}{@{}c|c|c|c@{}}[corners]
							\left( \frac{(x_j-x_i)(1+(1+w) x_i x_j)}{(x_i+x_j+w x_i x_j)(1-x_i x_j)} \right)_{1 \le i < j \le n} & \Block{2-1}{\left( 1 \right)_{1 \le i \le n}} & \Block{2-1}{ \left( \frac{1-x_i}{x_i} \left( \prod\limits_{\myatop{1 \le p \le n}{p \neq i}} \frac{1+ w x_i + x_i x_p}{x_i + x_p + w x_i x_p} \right) \right)_{1 \le i \le n}}  & \Block{2-1}{(0)_{1 \le i \le n}} \\ \cline{1}
							& & & \\ \cline{2-4}
							& & 0 & 1 \\ \cline{3-4}
							& & & \prod_{k=1}^{n} \frac{1}{x_k}
						\end{NiceArray}
					\end{pmatrix}$}}\\
			= \pf_{1 \le i < j \le n+1} {\scalebox{\myscale}{$\begin{pmatrix}
						\begin{NiceArray}{@{}c|c@{}}[corners]
							\left( \frac{(x_j-x_i)(1+(1+w) x_i x_j)}{(x_i+x_j+w x_i x_j)(1-x_i x_j)} \right)_{1 \le i < j \le n} & \Block{2-1}{\left(1 \right)_{1 \le i \le n}} \\ \cline{1}
							&
						\end{NiceArray}
					\end{pmatrix}$}}.
		\end{multline}
		To this end, we perform similar elementary row and column operations as in the even case on the matrix underlying the Pfaffian in the left-hand side of the previous equation. We multiply the $i$th column by
		\begin{equation*}
			-\frac{1+x_i}{x_i} \prod_{\myatop{1 \le p \le n}{p \neq i}} \frac{1-x_i x_p}{x_i -x_p}
		\end{equation*}
		and add it to the penultimate column for every $i\in\{1,\dots,n\}$; proceed analogously for the rows. We then obtain for the left-hand side of~\eqref{eq:2Pfaffians_Odd} the following due to Lemma~\ref{lem:MatrixOperationEven}:	
		\begin{equation}\label{eq:Pfaffian_odd}
			\pf_{1 \le i < j \le n+3}{\scalebox{\myscale}{$\begin{pmatrix}
						\begin{NiceArray}{@{}c|c|c|c@{}}[corners]
							\left( \frac{(x_j-x_i)(1+(1+w) x_i x_j)}{(x_i+x_j+w x_i x_j)(1-x_i x_j)} \right)_{1 \le i < j \le n} & \Block{2-1}{\left( 1 \right)_{1 \le i \le n}} & \Block{2-1}{\left(-1+\prod_{k=1}^{n} \frac{1}{x_k} \right)_{1 \le i \le n}}  & \Block{2-1}{(0)_{1 \le i \le n}} \\ \cline{1}
							& & & \\ \cline{2-4}
							& & 1+\prod_{k=1}^{n} \frac{1}{x_k} & 1 \\ \cline{3-4}
							& & & \prod_{k=1}^{n} \frac{1}{x_k}
						\end{NiceArray}
					\end{pmatrix}$}}.
		\end{equation}
		We continue with further row and column operations regarding~\eqref{eq:Pfaffian_odd}: we add $1-\prod_{k=1}^{n}\frac{1}{x_k}$ times the $(n+1)$th column and $-(1+\prod_{k=1}^{n}\frac{1}{x_k})$ times the $(n+3)$th column to the $(n+2)$th column; similarly for the rows. Finally, we subtract the $(n+1)$ row and column from the $(n+2)$th row and column, respectively. We eventually obtain	
		\begin{equation*}
			\pf_{1 \le i < j \le n+3}{\scalebox{\myscale}{$\begin{pmatrix}
						\begin{NiceArray}{@{}c|c|c|c@{}}[corners]
							\left( \frac{(x_j-x_i)(1+(1+w) x_i x_j)}{(x_i+x_j+w x_i x_j)(1-x_i x_j)} \right)_{1 \le i < j \le n} & \Block{2-1}{\left( 1 \right)_{1 \le i \le n}} & \Block{3-1}{\left(0 \right)_{1 \le i \le n+1}}  & \Block{3-1}{(0)_{1 \le i \le n+1}} \\ \cline{1}
							& & & \\ \cline{2-2}
							& & &  \\ \cline{3-4}
							& & & 1
						\end{NiceArray}
					\end{pmatrix}$}}.
		\end{equation*}
		Expanded along the last column, this expression finally coincides with the right-hand side of~\eqref{eq:2Pfaffians_Odd}. 	
	\end{proof}

\begin{thebibliography}{43}
\bibitem[1]{And79} Andrews, George E. Plane partitions. III. The weak Macdonald conjecture, Invent. Math., Volume 53 (1979) no. 3, pp. 193-225
\bibitem[2]{And94} Andrews, George E. Plane partitions. V. The TSSCPP conjecture, J. Combin. Theory Ser. A, Volume 66 (1994) no. 1, pp. 28-39
\bibitem[3]{ABF20} Ayyer, Arvind; Behrend, Roger E.; Fischer, Ilse Extreme diagonally and antidiagonally symmetric alternating sign matrices of odd order, Adv. Math., Volume 367 (2020), Paper no. 107125, 56 pages
\bibitem[4]{BFK23} Behrend, Roger E.; Fischer, Ilse; Koutschan, Christoph Diagonally symmetric alternating sign matrices, 2023, \url{https://arxiv.org/abs/2309.08446}
\bibitem[5]{BW16} Betea, D.; Wheeler, M. Refined Cauchy and Littlewood identities, plane partitions and symmetry classes of alternating sign matrices, J. Combin. Theory Ser. A, Volume 137 (2016), pp. 126-165
\bibitem[6]{BWZJ15} Betea, D.; Wheeler, M.; Zinn-Justin, P. Refined Cauchy/Littlewood identities and six-vertex model partition functions: II. Proofs and new conjectures, J. Algebraic Combin., Volume 42 (2015) no. 2, pp. 555-603
\bibitem[7]{BP18} Borodin, Alexei; Petrov, Leonid Higher spin six vertex model and symmetric rational functions, Selecta Math. (N.S.), Volume 24 (2018) no. 2, pp. 751-874
\bibitem[8]{BP99} Bressoud, David; Propp, James How the alternating sign matrix conjecture was solved, Notices Amer. Math. Soc., Volume 46 (1999) no. 6, pp. 637-646
\bibitem[9]{Bre00} Bressoud, David M. Elementary proofs of identities for Schur functions and plane partitions, Ramanujan J., Volume 4 (2000) no. 1, pp. 69-80
\bibitem[10]{Ciucu97} Ciucu, Mihai Enumeration of perfect matchings in graphs with reflective symmetry, J. Combin. Theory Ser. A, Volume 77 (1997) no. 1, pp. 67-97
\bibitem[11]{Elk92a} Elkies, Noam; Kuperberg, Greg; Larsen, Michael; Propp, James Alternating-sign matrices and domino tilings. I, J. Algebraic Combin., Volume 1 (1992) no. 2, pp. 111-132
\bibitem[12]{Elk92b} Elkies, Noam; Kuperberg, Greg; Larsen, Michael; Propp, James Alternating-sign matrices and domino tilings. II, J. Algebraic Combin., Volume 1 (1992) no. 3, pp. 219-234
\bibitem[13]{Fis07} Fischer, Ilse A new proof of the refined alternating sign matrix theorem, J. Combin. Theory Ser. A, Volume 114 (2007) no. 2, pp. 253-264
\bibitem[14]{Fis19Trap} Fischer, Ilse A constant term approach to enumerating alternating sign trapezoids, Adv. Math., Volume 356 (2019), Paper no. 106792, 23 pages
\bibitem[15]{Fis19} Fischer, Ilse Enumeration of alternating sign triangles using a constant term approach, Trans. Amer. Math. Soc., Volume 372 (2019) no. 2, pp. 1485-1508
\bibitem[16]{Fis23} Fischer, Ilse Computations versus bijections for tiling enumeration, Adv. in Appl. Math., Volume 142 (2023), Paper no. 102427, 49 pages
\bibitem[17]{Fis22} Fischer, Ilse Bounded Littlewood identity related to alternating sign matrices, Forum Math. Sigma, Volume 12 (2024), Paper no. e124, 44 pages
\bibitem[18]{FK20} Fischer, Ilse; Konvalinka, Matjaž The mysterious story of square ice, piles of cubes, and bijections, Proc. Natl. Acad. Sci. USA, Volume 117 (2020) no. 38, pp. 23460-23466
\bibitem[19]{FisRi2013} Fischer, Ilse; Riegler, Lukas Combinatorial reciprocity for monotone triangles, J. Combin. Theory Ser. A, Volume 120 (2013) no. 7, pp. 1372-1393
\bibitem[20]{FS23} Fischer, Ilse; Schreier-Aigner, Florian The relation between alternating sign matrices and descending plane partitions: $n + 3$ pairs of equivalent statistics, Adv. Math., Volume 413 (2023), p. Paper No. 108831, 47
\bibitem[21]{gangl} Gangl, Moritz Alternating sign pentagons and Magog pentagons, Adv. Math., Volume 474 (2025), Paper no. 110315, 29 pages
\bibitem[22]{Gav23} Gavrilova, Svetlana Refined Littlewood identity for spin Hall-Littlewood symmetric rational functions, Algebr. Comb., Volume 6 (2023) no. 1, pp. 37-51
\bibitem[23]{Hon21} Höngesberg, Hans Refined enumeration of halved monotone triangles and applications to vertically symmetric alternating sign trapezoids, J. Combin. Theory Ser. A, Volume 177 (2021), Paper no. 105336, 37 pages
\bibitem[24]{Hon22} Höngesberg, Hans A fourfold refined enumeration of alternating sign trapezoids, Electron. J. Combin., Volume 29 (2022) no. 3, Paper no. 3.42, 27 pages
\bibitem[25]{IN13} Ikeda, Takeshi; Naruse, Hiroshi ${K}$-theoretic analogues of factorial Schur ${P}$- and ${Q}$-functions, Adv. Math., Volume 243 (2013), pp. 22-66
\bibitem[26]{Kaw91} Kawanaka, Noriaki On subfield symmetric spaces over a finite field, Osaka J. Math., Volume 28 (1991) no. 4, pp. 759-791
\bibitem[27]{Kaw99} Kawanaka, Noriaki A $q$-series identity involving Schur functions and related topics, Osaka J. Math., Volume 36 (1999) no. 1, pp. 157-176
\bibitem[28]{Kra16} Krattenthaler, C. Plane partitions in the work of Richard Stanley and his school, The mathematical legacy of Richard P. Stanley, Amer. Math. Soc., Providence, RI, 2016, pp. 231-261
\bibitem[29]{Kra96} Krattenthaler, Christian A Gog-Magog conjecture, 1996 https://www.mat.univie.ac.at/\textasciitilde{}kratt/artikel/magog.html (unpublished note)
\bibitem[30]{Kup96} Kuperberg, Greg Another proof of the alternating-sign matrix conjecture, Internat. Math. Res. Notices (1996) no. 3, pp. 139-150
\bibitem[31]{Lit50} Littlewood, Dudley E. The Theory of Group Characters and Matrix Representations of Groups, Oxford University Press, New York, 1940, viii+292 pages
\bibitem[32]{MRR83} Mills, W. H.; Robbins, David P.; Rumsey, Howard Jr. Alternating sign matrices and descending plane partitions, J. Combin. Theory Ser. A, Volume 34 (1983) no. 3, pp. 340-359
\bibitem[33]{MRR86} Mills, W. H.; Robbins, David P.; Rumsey, Howard Jr. Self-complementary totally symmetric plane partitions, J. Combin. Theory Ser. A, Volume 42 (1986) no. 2, pp. 277-292
\bibitem[34]{P21} Petrov, Leonid Refined Cauchy identity for spin Hall-Littlewood symmetric rational functions, J. Combin. Theory Ser. A, Volume 184 (2021), Paper no. 105519, 50 pages
\bibitem[35]{Propp23} Propp, James Generalized domino-shuffling, Theoret. Comput. Sci., Volume 303 (2003) no. 2-3, pp. 267-301
\bibitem[36]{Rob91} Robbins, David P. The story of $1,2,7,42,429,7436,\cdots $, Math. Intelligencer, Volume 13 (1991) no. 2, pp. 12-19
\bibitem[37]{Rob} Robbins, David P. Symmetry classes of alternating sign matrices, 2000, \url{https://arxiv.org/abs/math/0008045}
\bibitem[38]{RR86} Robbins, David P.; Rumsey, Howard Jr. Determinants and alternating sign matrices, Adv. in Math., Volume 62 (1986) no. 2, pp. 169-184
\bibitem[39]{Sch11} Schur, J. Über die Darstellung der symmetrischen und der alternierenden Gruppe durch gebrochene lineare Substitutionen, J. Reine Angew. Math., Volume 139 (1911), pp. 155-250
\bibitem[40]{Sta99} Stanley, Richard P. Enumerative combinatorics. Vol. 2, Cambridge Studies in Advanced Mathematics, 62, Cambridge University Press, Cambridge, 1999, xii+581 pages
\bibitem[41]{War08} Warnaar, S. Ole Bisymmetric functions, Macdonald polynomials and $\mathfrak{sl}_3$ basic hypergeometric series, Compos. Math., Volume 144 (2008) no. 2, pp. 271-303
\bibitem[42]{WZJ16} Wheeler, Michael; Zinn-Justin, Paul Refined Cauchy/Littlewood identities and six-vertex model partition functions: III. Deformed bosons, Adv. Math., Volume 299 (2016), pp. 543-600
\bibitem[43]{Zei96} Zeilberger, Doron Proof of the alternating sign matrix conjecture, Electron. J. Combin., Volume 3 (1996) no. 2, Paper no. 13, 84 pages
\end{thebibliography}
\end{document}
